\subsection*{Other proofs of the fundamental theorems}

In {[}8{]} F. Schur showed that in canonical coordinates the \(\psi_{ik}\) of (15) satisfy the differential equations

\[
\frac {d}{d t} \left(t \psi_ {i k} (t a)\right) = \delta_ {i k} + \sum_ {j, l} c _ {j l} ^ {k} t a _ {l} \psi_ {i j} (t a).\tag{23}
\]

These are integrated and give a formula equivalent to the formula

\[
\varpi (\mathrm{X}) = \sum_ {n > 0} \frac {1}{(n + 1) !} (\mathrm{ad} (\mathrm{X})) ^ {n}\tag{24}
\]

of our Chapter III, §6, no. 4, Proposition 12; in particular, in canonical coordinates, the \(\psi_{ij}\) can be extended to integral functions of the \(a_k\) . F. Schur deduced a result which made precise an earlier remark of Lie: if, in definition (4) of transformation groups, it is only assumed that the \(f_i\) are of class \(\mathbf{C}^2\) , then the group is holohedrically isomorphic to an analytic group. \(\dagger\) Following his research on the integration of differential systems, E. Cartan ({[}12{]}, vol.~II \(_{2}\) , p.~371) introduced in 1904 Pfaffian forms

\[
\omega_ {k} = \sum_ {i = 1} ^ {r} \psi_ {k i} d a _ {i} (1 \leqslant i \leqslant r)\tag{25}
\]

(in the notation of (15)), called later Maurer-Cartan forms. The Maurer conditions (17) could be written

\[
d \omega_ {k} = - \frac {1}{2} \sum_ {i, j} c _ {i j} ^ {k} \omega_ {i} \wedge \omega_ {j};
\]

E. Cartan showed that the theory of finite continuous groups could be developed starting from the \(\omega_{k}\) and established the equivalence of this point of view and that of Lie. But, for him, the interest of this method was above all that it could be adapted to ``infinite continuous groups'' the theory of which he pushed much further than Lie had done and which allowed him to develop his theory of the generalized ``moving frame''.

